\documentclass[../Main.tex]{subfiles}

\begin{document}

\section{Reflexion} \label{sec:Reflexion}

\subsection{Reflexion des Arbeitsprozesses}

\begin{itemize}[leftmargin=1.5cm]
    \item \textcolor{red}{\textbf{1. Ausgangslage und Plan:} Womit seid ihr gestartet (ursprüngliche Projektidee)? Was war laut Vertrag geplant? Meilensteine: Zwischenprodukt 10.6., Zwischenpräsentation 17.6., Analyse der Referenzspiele Anfang Juli, Spielsysteme Mitte August, Doku und Fragebögen Ende September, Playtests Mitte Oktober, Verbesserungen Ende Oktober.}
    \item \textcolor{red}{\textbf{2. Scope-Entscheide:} Am 20.5. wurde der Projektumfang reduziert -- was wurde gestrichen und warum? Später: Bonus-Inhalte weggelassen (z.\,B. Boss, Lock-on), technische Dokumentation ausgelagert, README gestrichen, nicht alle Recherchepunkte aus Kapitel~\ref{sec:Leveldesign} umgesetzt.}
    \item \textcolor{red}{\textbf{3. Plan und tatsächlicher Verlauf:} Wo seid ihr im Plan geblieben, wo nicht (z.\,B. Playtests 26.10.--2.11. statt Mitte Oktober)? Warum?}
    \item \textcolor{red}{\textbf{4. Was gut funktioniert hat:} z.\,B. iteratives Vorgehen, Versionsverwaltung mit Git, Playtest-Vorgehen -- jeweils mit einem Grund.}
    \item \textcolor{red}{\textbf{5. Schwierigkeiten und wie ihr sie gelöst habt:} z.\,B. Zeitmanagement, Umsetzung der Rückmeldungen aus der Zwischenbewertung.}
    \item \textcolor{red}{\textbf{6. Methodik im Rückblick:} Haben sich MDA-Fragebogen, Leveldesign-Recherche und Playtest-Fragebögen bewährt? Kurz, die Einzelheiten stehen in den Limitationen (Abschnitt~\ref{sec:Limitationen}).}
    \item \textcolor{red}{\textbf{7. Was ihr anders machen würdet} und was ihr daraus für künftige Projekte mitnehmt.}
    \item \textcolor{red}{\textbf{8. Verweis aufs Arbeitsjournal}, wo der Verlauf im Detail dokumentiert ist.}
\end{itemize}

\subsection{Gruppenarbeit und Aufgabenteilung}

\begin{itemize}[leftmargin=1.5cm]
    \item \textcolor{red}{\textbf{1. Wie ihr die Arbeit aufgeteilt habt:} Die Zwischenbewertung nennt: Rodin stärker bei Animation, UI, Bericht und README; Zhile stärker bei Architektur, Bewegung, State Machines und Kampf-Grundlagen.}
    \item \textcolor{red}{\textbf{2. Tabelle Aufgabenteilung und Eigenleistung} für Bericht und Prototyp (Rodin / Zhile / gemeinsam) -- ausdrückliche Forderung der Zwischenbewertung.}
    \item \textcolor{red}{\textbf{3. Zusammenarbeit und Kommunikation:} Wie habt ihr euch abgesprochen (Treffen, Discord, GitHub-Issues \ldots)? Wie habt ihr Entscheidungen getroffen?}
    \item \textcolor{red}{\textbf{4. Veränderung im Verlauf:} Hat sich die Aufteilung verändert, z.\,B. in der Schlussphase? Warum?}
    \item \textcolor{red}{\textbf{5. Fazit:} Was lief gut, was war schwierig, was würdet ihr in einer Gruppenarbeit anders machen?}
\end{itemize}

\subsection{Verwendung von KI} \label{sec:ReflexionKI}

\begin{itemize}[leftmargin=1.5cm]
    \item \textcolor{red}{\textbf{1. Welche Tools wofür:} kurz zusammengefasst, Details in der KI-Deklaration (Abschnitt~\ref{sec:KIDeklaration}).}
    \item \textcolor{red}{\textbf{2. Was geholfen hat:} z.\,B. Sprachkorrektur, Struktur, Rückmeldungen zu Widersprüchen, Fehlersuche im Code.}
    \item \textcolor{red}{\textbf{3. Grenzen und Probleme:} Wo lag die KI falsch oder hat Dinge angenommen, die nicht stimmten? Wo musstet ihr selbst entscheiden?}
    \item \textcolor{red}{\textbf{4. Wie ihr die Ergebnisse überprüft habt:} z.\,B. Abgleich mit Quellen und Playtest-Protokollen, eigene Kontrolle vor der Übernahme.}
    \item \textcolor{red}{\textbf{5. Eigenanteil sichern:} Wie habt ihr darauf geachtet, dass Inhalte, Interpretation und Schlussfolgerungen von euch stammen (z.\,B. KI-formulierte Stellen in eigenen Worten neu geschrieben)?}
    \item \textcolor{red}{\textbf{6. Fazit:} Würdet ihr KI wieder so einsetzen? Was würdet ihr anders machen?}
\end{itemize}

\subsection{Persönliche Erkenntnisse}

\subsubsection{Rodin Arslan}

\begin{itemize}[leftmargin=1.5cm]
    \item \textcolor{red}{\textbf{1. Fachlich:} Was hast du gelernt (z.\,B. Godot, Animation, UI, Leveldesign, Playtesting, \LaTeX)?}
    \item \textcolor{red}{\textbf{2. Methodisch:} z.\,B. Recherche und Zitieren, wissenschaftliches Schreiben, Projektplanung.}
    \item \textcolor{red}{\textbf{3. Persönlich:} z.\,B. Zeitmanagement, Durchhalten, Umgang mit Rückmeldungen.}
    \item \textcolor{red}{\textbf{4. Was du mitnimmst} für die Zukunft.}
\end{itemize}

\subsubsection{Zhile Zhu}

\textcolor{red}{ZHILE DAS IST FÜR DICH}

\end{document}
